% Chapter 3 abstract (edit this file in the chapter Subfiles folder;
% included from Chp3.tex and from the JHU dissertation 08-chapter-3.tex).
The inequality literature is notable in its ability to bring together four distinct concepts: (i) measures of inequality, (ii) social welfare functions, (iii) mean values, and (iv) models of choice under uncertainty, as noted by \cite{hd20}. This interplay between statistical and theoretical analysis has been commonly applied in the context of inequality in the distribution of \textit{income}. It is well-known that income and wealth inequality are markedly distinct phenomena, and this distinction becomes more pronounced when racial considerations are central. The literature has long documented that wealth is more strongly skewed than income, especially in the upper tail \parencite{jbab18}. Notably, as recent as 2019, it has been reported that \say{the median white household held \$188,200 in wealth -- 7.8 times that of the typical Black household}. This paper provides a theoretical foundation for wealth inequality that incorporates the first three of the four distinct concepts. The ranking over wealth distributions is established through a choice under ambiguity model, and the \textit{social welfare approach} to inequality measures in \cite{aa71} is applied in this setting to derive a measure of wealth inequality aversion. Using the triennial data on wealth in the Survey of Consumer Finances (SCF) from 1989--2022, I show that the choice under objective uncertainty inequality measure underestimates wealth inequality compared to the measure proposed via choice under ambiguity. Both theoretical measures of inequality aversion track the general trend in the Black-White median wealth gap over the same time period.
